\documentclass{article}%
\usepackage[T1]{fontenc}%
\usepackage[utf8]{inputenc}%
\usepackage{lmodern}%
\usepackage{textcomp}%
\usepackage{lastpage}%
\usepackage{geometry}%
\geometry{margin=1in,headheight=14pt,headsep=25pt}%
\usepackage{hyperref}%
\usepackage{enumitem}%
\usepackage{fancyhdr}%
\usepackage{xcolor}%
\usepackage{url}%
\usepackage{breakurl}%
%

            \hypersetup{
                colorlinks=true,
                linkcolor=blue,
                filecolor=magenta,
                urlcolor=blue
            }
            \pagestyle{fancy}
            \fancyhf{}
            \rhead{Generated on \today}
            \cfoot{\thepage}
            
            % Configure enumeration settings
            \setlist[enumerate,1]{label=\arabic*.}
            \setlist[enumerate,2]{label=\alph*.}
            \setlist[enumerate,3]{label=\Alph*.}
            \setlist[enumerate]{itemsep=0.5em}
            
            % Define a command for red text
            \newcommand{\incorrect}[1]{\textcolor{red}{#1}}
        %
\lhead{2024{-}11{-}05{-}NetFlows}%
\title{Summary for 2024{-}11{-}05{-}NetFlows}%
\author{Generated by Scribe.AI}%
\date{\today}%
%
\begin{document}%
\normalsize%
\maketitle%
\section{Summary}%
\label{sec:Summary}%
This lecture on Network Flow within a Linear Programming course explored the formulation and solution of minimum-cost network flow problems.

* **Network:** A network is defined as a tuple (N, A), where N is the set of nodes and A is the set of directed arcs connecting those nodes.

* **Network Flow Problem:**  The minimum-cost network flow problem aims to minimize the total cost of transporting flow across a network, subject to supply and demand constraints at each node.

* **Flow Conservation:** At each node, the total inflow plus the node's supply must equal the total outflow.  This constraint is represented mathematically by the equation Ax = -b, where A is the network's incidence matrix, x is the flow vector, and b is the supply/demand vector.

* **Incidence Matrix:** The incidence matrix A represents the network's structure, with each column corresponding to an arc and each row to a node.  A non-zero entry indicates the arc's incidence with the node.

* **Linear Programming Formulation:** The minimum-cost network flow problem is formulated as a linear program, minimizing c<sup>T</sup>x subject to Ax = -b and x ≥ 0, where c is the vector of arc costs.

* **Duality:** The dual linear program provides valuable insights into the problem's solution, with dual variables representing node potentials and dual slack variables representing the cost differences between potential and actual arc costs.

* **Spanning Trees and Bases:**  Spanning trees in the network correspond to bases in the simplex method used to solve the linear program.  A basic feasible solution is characterized by zero flow on arcs not in the spanning tree.

* **Optimality Conditions:** A basic feasible solution is optimal if all flows are non-negative and all dual slacks are non-negative.  The primal simplex method can be used to iteratively improve a primal feasible but not dual feasible solution.

%
\end{document}